\section{作文模板}

\subsection{书信类}

\subsubsection{建议信}

\begin{center}
\textbf{Dear \_\_\_\_\_\_ ,}
\end{center}

I'm sorry to hear that you are having trouble with \_\_\_\_\_\_ . I'm writing to offer you some practical suggestions.

First and foremost, I suggest that you \_\_\_\_\_\_ (建议1). This will help you \_\_\_\_\_\_ (好处). Additionally, it would be a good idea to \_\_\_\_\_\_ (建议2). By doing so, you can \_\_\_\_\_\_ (好处). Last but not least, why not \_\_\_\_\_\_ (建议3)? This approach has proved effective for many people.

I hope these suggestions will be of some help to you. If there is anything else I can do, please don't hesitate to let me know.

\begin{center}
Yours sincerely,\\
Li Hua
\end{center}

\subsubsection{邀请信}

\begin{center}
\textbf{Dear \_\_\_\_\_\_ ,}
\end{center}

I am writing to invite you to \_\_\_\_\_\_ (活动名称), which will be held on \_\_\_\_\_\_ (日期) at \_\_\_\_\_\_ (地点).

The activity will begin at \_\_\_\_\_\_ (开始时间) and last approximately \_\_\_\_\_\_ (持续时间). First, we will \_\_\_\_\_\_ (活动内容1). After that, there will be \_\_\_\_\_\_ (活动内容2). I believe it will be a wonderful experience and you will enjoy it.

I would be honored if you could come. Please let me know if you are available.

\begin{center}
Yours,\\
Li Hua
\end{center}

\subsubsection{感谢信}

\begin{center}
\textbf{Dear \_\_\_\_\_\_ ,}
\end{center}

I am writing to express my heartfelt gratitude for \_\_\_\_\_\_ . Without your help, I couldn't have \_\_\_\_\_\_ (如果没有帮助会怎样).

Your \_\_\_\_\_\_ (具体的帮助) made a great difference to me. It not only \_\_\_\_\_\_ (帮助1) but also \_\_\_\_\_\_ (帮助2). I am truly grateful for your kindness and generosity.

I hope to have the opportunity to return your favor in the future. Thank you again for everything.

\begin{center}
Yours sincerely,\\
Li Hua
\end{center}

\subsubsection{求助信}

\begin{center}
\textbf{Dear \_\_\_\_\_\_ ,}
\end{center}

I hope this message finds you well. I am writing to ask for your help with \_\_\_\_\_\_ .

The problem is that \_\_\_\_\_\_ (困难描述). I have tried \_\_\_\_\_\_ (已尝试的方法), but it didn't work. I would be extremely grateful if you could \_\_\_\_\_\_ (求助内容). Specifically, I need advice on \_\_\_\_\_\_ (具体方面).

I apologize for troubling you, but I really need your assistance. Thank you in advance for your kind help.

\begin{center}
Yours,\\
Li Hua
\end{center}

\subsection{通知类}

\subsubsection{书面通知}

\begin{center}
\textbf{Notice}
\end{center}

Dear students,

In order to \_\_\_\_\_\_ (目的), our school is going to organize \_\_\_\_\_\_ (活动). The details are as follows.

The activity will take place in \_\_\_\_\_\_ (地点) from \_\_\_\_\_\_ (时间) to \_\_\_\_\_\_ (时间) on \_\_\_\_\_\_ (日期). All students are welcome to participate. During the activity, we will \_\_\_\_\_\_ (活动内容).

Those who are interested are expected to \_\_\_\_\_\_ (报名方式). Please sign up before \_\_\_\_\_\_ (截止日期). For more information, please contact us at \_\_\_\_\_\_ (联系方式).

We are looking forward to your participation.

\begin{center}
The Students' Union\\
\_\_\_\_\_\_ (日期)
\end{center}

\subsubsection{口头通知}

\begin{center}
\textbf{Announcement}
\end{center}

May I have your attention, please? I have an announcement to make.

\_\_\_\_\_\_ (活动) will be held \_\_\_\_\_\_ (时间/地点). Everyone is required to \_\_\_\_\_\_ (要求). Please \_\_\_\_\_\_ (注意事项).

That's all. Thank you.

\subsection{读后续写}

\subsubsection{情绪描写素材}

\begin{center}
\begin{tabular}{c|l}
\textbf{情绪} & \textbf{描写句} \\ \toprule
高兴 & Her face lit up with a radiant smile. \\
 & A wave of joy washed over her. \\
 & His eyes sparkled with happiness. \\
悲伤 & Tears welled up in her eyes. \\
 & His heart felt heavy with sorrow. \\
 & She choked back a sob. \\
紧张 & His heart was pounding like a drum. \\
 & She bit her lip, her hands trembling. \\
 & He paced back and forth, unable to sit still. \\
惊讶 & Her jaw dropped in disbelief. \\
 & He stood frozen, eyes wide with shock. \\
 & She couldn't believe what she was seeing. \\
恐惧 & A cold shiver ran down his spine. \\
 & Her face turned pale with terror. \\
 & He felt a knot in his stomach. \\
\bottomrule
\end{tabular}
\end{center}

\subsubsection{动作描写素材}

\begin{center}
\begin{tabular}{c|l}
\textbf{动作类别} & \textbf{常用表达} \\ \toprule
看 & glance at, stare at, glare at, peek at, catch sight of \\
走 & march, wander, stagger, tiptoe, dash, rush, stride \\
说 & whisper, murmur, announce, exclaim, reply, shout \\
手部 & grasp, clutch, squeeze, reach out, wave, grab \\
坐/站 & sink into, lean against, rise to one's feet \\
跑/跳 & dash off, sprint, leap, jump to one's feet \\
\bottomrule
\end{tabular}
\end{center}

\subsubsection{环境描写素材}

\begin{itemize}
  \item \textbf{晴天}：The sun was shining brightly in the clear blue sky.
  \item \textbf{阴天}：Dark clouds gathered overhead, threatening rain.
  \item \textbf{雨天}：The rain poured down like a curtain.
  \item \textbf{夜晚}：The stars twinkled in the velvet sky.
  \item \textbf{森林}：Tall trees stood silently, their leaves rustling in the breeze.
  \item \textbf{海边}：Waves crashed against the shore, spraying foam into the air.
\end{itemize}

\subsubsection{续写万能结尾}

\begin{itemize}
  \item 感悟型：This experience taught me that \_\_\_\_\_\_ .
  \item 成长型：From then on, I \_\_\_\_\_\_ and never looked back.
  \item 感恩型：I will always remember \_\_\_\_\_\_ for what he/she did for me.
  \item 总结型：Sometimes, the smallest acts of kindness can make the biggest difference.
\end{itemize}

\subsection{议论文}

\subsubsection{观点对比类}

\begin{center}
\textbf{Should We \_\_\_\_\_\_ ?}
\end{center}

Recently, the topic of \_\_\_\_\_\_ has sparked a heated debate. Different people hold different views.

On one hand, some people believe that \_\_\_\_\_\_ (观点A). They argue that \_\_\_\_\_\_ (理由1). Furthermore, \_\_\_\_\_\_ (理由2).

On the other hand, others insist that \_\_\_\_\_\_ (观点B). According to them, \_\_\_\_\_\_ (理由1). Besides, \_\_\_\_\_\_ (理由2).

In my opinion, \_\_\_\_\_\_ (我的观点). It is high time that we \_\_\_\_\_\_ (建议).

\subsubsection{现象分析类}

\begin{center}
\textbf{\_\_\_\_\_\_ : A Growing Trend}
\end{center}

In recent years, \_\_\_\_\_\_ has become increasingly common. There are several reasons behind this phenomenon.

First and foremost, \_\_\_\_\_\_ (原因1). Moreover, \_\_\_\_\_\_ (原因2). Last but not least, \_\_\_\_\_\_ (原因3).

As far as I am concerned, \_\_\_\_\_\_ (个人看法). On the one hand, \_\_\_\_\_\_ (积极方面). On the other hand, \_\_\_\_\_\_ (消极方面). Therefore, it is necessary that we \_\_\_\_\_\_ (对策).

\subsection{高分替换词汇}

\begin{center}
\begin{tabular}{c|c}
\textbf{普通词} & \textbf{高分替换词} \\ \toprule
important & significant, essential, crucial, vital \\
many & numerous, a variety of, a large number of \\
good & excellent, outstanding, remarkable, fantastic \\
bad & terrible, awful, negative, harmful \\
beautiful & gorgeous, splendid, magnificent \\
big & enormous, massive, considerable \\
small & tiny, minor, slight \\
like & be fond of, be keen on, have a passion for \\
think & believe, hold the view that, be convinced that \\
because & due to, owing to, on account of, in that \\
so & therefore, consequently, as a result, thus \\
but & however, nevertheless, whereas, yet \\
very & extremely, thoroughly, deeply, highly \\
\bottomrule
\end{tabular}
\end{center}

\subsection{写作检查清单}

\subsubsection{写前}

\begin{itemize}
  \item 审题：文体 / 人称 / 时态 / 要点全覆盖
  \item 列提纲：开头 $\to$ 正文 $\to$ 结尾
\end{itemize}

\subsubsection{写中}

\begin{itemize}
  \item 注意语法正确（时态、主谓一致、名词单复数）
  \item 句式多样化（复合句、倒装句、非谓语）
  \item 词汇多样化（使用替换词）
  \item 连接词自然（first, moreover, however, therefore, in conclusion）
\end{itemize}

\subsubsection{写后检查}

\begin{itemize}
  \item 是否有拼写错误
  \item 是否有要点遗漏
  \item 字数是否达标
  \item 书写是否工整
\end{itemize}
